% =====================================================================
% Chủ đề 5 — Tư duy logic và Câu đố — phần c: bài 31–40 (mở rộng)
% =====================================================================
% ---- helper macros for parts c/d of topic 5 (prefix \lgxcd, guarded) ----
% shape codes: 1 = circle, 2 = square, 3 = triangle, 4 = star
\providecommand{\lgxcdShape}[3]{\ifcase#1\or
  \draw[Ink,fill=Slate!35,line width=0.6pt] (#2,#3) circle (0.19);\or
  \draw[Ink,fill=tmBlue!35,line width=0.6pt] ($(#2,#3)+(-0.18,-0.18)$) rectangle ($(#2,#3)+(0.18,0.18)$);\or
  \draw[Ink,fill=tmOrange!45,line width=0.6pt] ($(#2,#3)+(-0.2,0.17)$)--($(#2,#3)+(0.2,0.17)$)--($(#2,#3)+(0,-0.19)$)--cycle;\or
  \draw[Ink,fill=tmGreen!45,line width=0.6pt,shift={(#2,#3)}] (90:0.24) \foreach \lgxcdk in {1,...,9}{ -- ({90+36*\lgxcdk}:{ifthenelse(mod(\lgxcdk,2)==1,0.1,0.24)})} -- cycle;\fi}
% inline shape in text
\providecommand{\lgxcdI}[1]{\tikz[baseline=-0.7ex]{\lgxcdShape{#1}{0}{0}}}
% a row of n shapes centred at the origin: \lgxcdRow{n}{codes}
\providecommand{\lgxcdRow}[2]{\foreach \lgxcdt [count=\lgxcdi] in {#2}{%
  \pgfmathsetmacro\lgxcdxx{(\lgxcdi-(#1+1)/2)*0.48}\lgxcdShape{\lgxcdt}{\lgxcdxx}{0.21}}}
% equal-arm balance at height #1: \lgxcdBal{y}{n left}{codes left}{n right}{codes right}
\providecommand{\lgxcdBal}[5]{\begin{scope}[yshift=#1cm]
  \draw[Ink,line width=1pt] (-2.7,0)--(-0.3,0) (0.3,0)--(2.7,0);
  \draw[Ink,line width=0.9pt] (-1.5,0)--(-1.5,-0.55)--(1.5,-0.55)--(1.5,0);
  \draw[Ink,line width=0.9pt] (-0.18,-0.85)--(0,-0.58)--(0.18,-0.85);
  \begin{scope}[xshift=-1.5cm]\lgxcdRow{#2}{#3}\end{scope}
  \begin{scope}[xshift=1.5cm]\lgxcdRow{#4}{#5}\end{scope}\end{scope}}
% square grid n x n with box lines every #2 rows / #3 columns; #4 = given r/c/v, #5 = filled r/c/v (rows from top)
\providecommand{\lgxcdSud}[5]{\begin{scope}[x=0.62cm,y=0.62cm]
  \foreach \lgxcdi in {0,...,#1}{\draw[Slate!60,line width=0.5pt] (0,\lgxcdi)--(#1,\lgxcdi) (\lgxcdi,0)--(\lgxcdi,#1);}
  \foreach \lgxcdi in {0,#2,...,#1}{\draw[Ink,line width=1.4pt] (0,\lgxcdi)--(#1,\lgxcdi);}
  \foreach \lgxcdi in {0,#3,...,#1}{\draw[Ink,line width=1.4pt] (\lgxcdi,0)--(\lgxcdi,#1);}
  \foreach \lgxcdr/\lgxcdc/\lgxcdv in {#4}{\node[font=\sffamily\bfseries,text=Ink] at (\lgxcdc+0.5,#1-\lgxcdr-0.5){\lgxcdv};}
  \foreach \lgxcdr/\lgxcdc/\lgxcdv in {#5}{\node[font=\sffamily,text=tmBlue] at (\lgxcdc+0.5,#1-\lgxcdr-0.5){\lgxcdv};}
  \end{scope}}

\clearpage
\dividerbar{Mở rộng -- Nâng cao}

% ---------------------------------------------------------------- Bài 31
\begin{probbox}{31}{\relax}
Trên một hòn đảo, mỗi người dân hoặc là \emph{người trung thực} (luôn nói thật), hoặc là \emph{kẻ nói dối} (luôn nói dối).
\begin{enumerate}[label=\alph*),leftmargin=1.6em,itemsep=1pt,topsep=2pt]
  \item Em gặp hai người dân Khánh và Lâm. Khánh nói: ``Cả hai chúng tôi đều là kẻ nói dối.'' Hỏi mỗi người thuộc loại nào?
  \item Em gặp ba người dân An, Bình và Chi.\\
  An nói: ``Bình là kẻ nói dối.''\\
  Bình nói: ``Chi là kẻ nói dối.''\\
  Chi nói: ``An và Bình đều là kẻ nói dối.''\\
  Hỏi mỗi người thuộc loại nào?
\end{enumerate}
\end{probbox}
\probsol{
\textbf{Phân tích đề.} Mỗi người chỉ có hai khả năng. Cách làm chắc chắn nhất là \emph{giả sử} một khả năng, rồi xem các câu nói có
phù hợp với giả sử đó không. Nếu dẫn tới mâu thuẫn thì loại khả năng đó.
\begin{recallbox}Người trung thực nói câu nào thì câu đó \textbf{đúng}. Kẻ nói dối nói câu nào thì câu đó \textbf{sai}.
Như vậy, mỗi câu nói cho ta một ``phương trình'': \emph{câu nói đúng} $\Longleftrightarrow$ \emph{người nói là người trung thực}.\end{recallbox}
\par\smallskip\textbf{Câu a).}
\step{1}{Giả sử Khánh là người trung thực}
Khi đó câu của Khánh đúng, tức là Khánh là kẻ nói dối -- trái với giả sử. Vậy Khánh \textbf{không thể} là người trung thực.
\step{2}{Khánh là kẻ nói dối}
Câu ``Cả hai chúng tôi đều là kẻ nói dối'' là câu sai. Câu này sai khi \emph{ít nhất một} người không phải kẻ nói dối. Khánh đã
là kẻ nói dối, nên người còn lại -- Lâm -- phải là người trung thực.
\begin{warnbox}Phủ định của ``cả hai đều nói dối'' \textbf{không phải} là ``cả hai đều nói thật'', mà là ``có ít nhất một người
nói thật''.\end{warnbox}
\par\smallskip\textbf{Câu b).}
\step{3}{Giả sử An là người trung thực}
Câu của An đúng $\Rightarrow$ Bình là kẻ nói dối $\Rightarrow$ câu của Bình sai $\Rightarrow$ Chi là người trung thực
$\Rightarrow$ câu của Chi đúng $\Rightarrow$ An là kẻ nói dối. Trái với giả sử! Loại.
\step{4}{Vậy An là kẻ nói dối}
Câu của An sai $\Rightarrow$ Bình \emph{không} nói dối, tức Bình là người trung thực $\Rightarrow$ câu của Bình đúng
$\Rightarrow$ Chi là kẻ nói dối.
\step{5}{Kiểm tra câu của Chi}
Chi là kẻ nói dối nên câu ``An và Bình đều là kẻ nói dối'' phải sai. Thật vậy, câu đó sai vì Bình là người trung thực.
Không có mâu thuẫn nào.
\begin{center}
\begin{tabular}{|c|c|l|c|}
\hline
\rowcolor{tmBlue!10} Người & Loại & Câu nói & Câu nói đúng/sai\\ \hline
An & nói dối & Bình là kẻ nói dối & sai \\ \hline
Bình & trung thực & Chi là kẻ nói dối & đúng\\ \hline
Chi & nói dối & An và Bình đều nói dối & sai\\ \hline
\end{tabular}
\end{center}
\verify{Trong bảng, mỗi người trung thực nói câu đúng, mỗi kẻ nói dối nói câu sai. Ta cũng đã thấy khả năng ``An trung thực''
bị loại, nên đây là lời giải duy nhất.}
\ans{a) Khánh nói dối, Lâm trung thực;\quad b) An và Chi nói dối, Bình trung thực}
}

% ---------------------------------------------------------------- Bài 32
\begin{probbox}{32}{\relax}
Bạn Mai nói dối vào các ngày thứ Hai, thứ Ba, thứ Tư và nói thật vào những ngày còn lại trong tuần. Bạn Hùng nói dối vào các ngày
thứ Năm, thứ Sáu, thứ Bảy và nói thật vào những ngày còn lại. Một hôm, cả hai bạn đều nói: ``Hôm qua là ngày tôi nói dối.''
Hỏi hôm đó là thứ mấy?
\end{probbox}
\probsol{
\textbf{Phân tích đề.} Câu nói ``Hôm qua là ngày tôi nói dối'' có thể đúng hoặc sai, tuỳ theo ngày. Với từng bạn, ta tìm những ngày
bạn đó \emph{có thể} nói câu này, rồi lấy ngày chung của hai bạn.
\begin{recallbox}Vào ngày nói thật, câu nói phải đúng; vào ngày nói dối, câu nói phải sai.\end{recallbox}
\step{1}{Lập bảng cho Mai (nói dối T2, T3, T4)}
\begin{center}\small
\begin{tabular}{|l|c|c|c|c|c|c|c|}
\hline
\rowcolor{tmBlue!10} Hôm nay & T2 & T3 & T4 & T5 & T6 & T7 & CN\\ \hline
Hôm nay Mai & dối & dối & dối & thật & thật & thật & thật\\ \hline
Hôm qua Mai & thật (CN) & dối & dối & dối & thật & thật & thật\\ \hline
Câu nói đúng? & sai & đúng & đúng & đúng & sai & sai & sai\\ \hline
Mai nói được? & \textbf{có} & không & không & \textbf{có} & không & không & không\\ \hline
\end{tabular}
\end{center}
Mai nói được câu này khi: hôm nay nói dối và câu sai (thứ Hai), hoặc hôm nay nói thật và câu đúng (thứ Năm).
\step{2}{Lập bảng cho Hùng (nói dối T5, T6, T7)}
\begin{center}\small
\begin{tabular}{|l|c|c|c|c|c|c|c|}
\hline
\rowcolor{tmBlue!10} Hôm nay & T2 & T3 & T4 & T5 & T6 & T7 & CN\\ \hline
Hôm nay Hùng & thật & thật & thật & dối & dối & dối & thật\\ \hline
Hôm qua Hùng & thật & thật & thật & thật & dối & dối & dối\\ \hline
Câu nói đúng? & sai & sai & sai & sai & đúng & đúng & đúng\\ \hline
Hùng nói được? & không & không & không & \textbf{có} & không & không & \textbf{có}\\ \hline
\end{tabular}
\end{center}
Hùng nói được câu này vào thứ Năm hoặc Chủ nhật.
\step{3}{Ngày chung}
Mai: $\{$thứ Hai, thứ Năm$\}$; Hùng: $\{$thứ Năm, Chủ nhật$\}$. Ngày chung duy nhất là \textbf{thứ Năm}.
\begin{tipbox}Quy luật nhanh: người này chỉ nói được câu ``Hôm qua tôi nói dối'' vào ngày \emph{chuyển chế độ}
(hôm qua và hôm nay khác chế độ): ngày đầu tiên nói dối, hoặc ngày đầu tiên nói thật trở lại.\end{tipbox}
\verify{Thứ Năm: Mai nói thật, hôm qua (thứ Tư) Mai nói dối -- câu đúng, hợp lí. Hùng nói dối, hôm qua (thứ Tư) Hùng nói thật --
câu sai, hợp lí với việc Hùng đang nói dối.}
\ans{Thứ Năm}
}

% ---------------------------------------------------------------- Bài 33
\begin{probbox}{33}{\relax}
Sáu bạn Lan, Minh, Nam, Oanh, Phúc và Quân chạy thi. Không có hai bạn nào về đích cùng lúc. Biết rằng:
\begin{itemize}[leftmargin=1.4em,itemsep=0pt,topsep=2pt]
  \item Minh về đích sau Lan nhưng trước Nam;
  \item Oanh về đích sau Nam;
  \item Phúc về đích \emph{ngay sau} Lan;
  \item Quân không về nhất cũng không về bét, và Quân về trước Minh.
\end{itemize}
Hãy xếp thứ tự về đích của sáu bạn. Ai về thứ ba?
\end{probbox}
\probsol{
\textbf{Phân tích đề.} Ta có các điều kiện ``trước -- sau''. Cách hay là ghép chúng thành một \emph{chuỗi} rồi chèn dần những
bạn còn lại vào. Viết $X \to Y$ nghĩa là ``$X$ về trước $Y$''.
\step{1}{Ghép điều kiện 1 và 2}
Lan $\to$ Minh $\to$ Nam và Nam $\to$ Oanh, nên: \ Lan $\to$ Minh $\to$ Nam $\to$ Oanh.
\step{2}{Chèn Phúc}
Phúc về \emph{ngay sau} Lan, tức là không ai chen giữa Lan và Phúc: \ Lan, Phúc $\to$ Minh $\to$ Nam $\to$ Oanh.
\step{3}{Chèn Quân}
Quân về trước Minh, nên Quân đứng ở một trong các chỗ: trước Lan, hoặc giữa Phúc và Minh (không được chen giữa Lan và Phúc).
Nhưng Quân không về nhất, nên Quân không đứng trước Lan. Vậy:
\[ \text{Lan} \to \text{Phúc} \to \text{Quân} \to \text{Minh} \to \text{Nam} \to \text{Oanh}. \]
\begin{figC}
\foreach \n/\x in {Lan/0,Phúc/1,Quân/2,Minh/3,Nam/4,Oanh/5}{
  \node[draw=tmBlue,fill=tmBlue!8,rounded corners=3pt,minimum width=1.5cm,minimum height=0.7cm] at (\x*1.9,0){\n};}
\foreach \x [count=\k] in {0,...,5}{\node[font=\scriptsize,text=Slate] at (\x*1.9,-0.6){hạng \k};}
\foreach \x in {0,...,4}{\draw[->,tmBlue,thick] (\x*1.9+0.78,0)--(\x*1.9+1.12,0);}
\end{figC}
\step{4}{Kiểm tra lại từng điều kiện}
Minh (hạng 4) sau Lan (hạng 1), trước Nam (hạng 5): đúng. Oanh (hạng 6) sau Nam: đúng. Phúc (hạng 2) ngay sau Lan: đúng.
Quân hạng 3 -- không nhất, không bét, trước Minh: đúng.
\begin{warnbox}Nếu quên chữ ``ngay sau'' thì Quân có thể chen giữa Lan và Phúc, và bài toán sẽ có hai đáp án. Hãy đọc kĩ từng chữ
của điều kiện!\end{warnbox}
\verify{Đã thử mọi vị trí có thể của Quân; chỉ có một thứ tự thoả mãn cả bốn điều kiện.}
\ans{Lan, Phúc, Quân, Minh, Nam, Oanh; Quân về thứ ba}
}

% ---------------------------------------------------------------- Bài 34
\begin{probbox}{34}{\relax}
Năm bạn Hà, Khoa, Linh, My và Tùng ngồi trên một hàng ghế gồm $5$ ghế đánh số $1, 2, 3, 4, 5$ từ trái sang phải. Biết rằng:
Hà và Khoa ngồi cạnh nhau; Linh và My \emph{không} ngồi cạnh nhau; Tùng ngồi ở một đầu hàng ghế (ghế $1$ hoặc ghế $5$).
\begin{enumerate}[label=\alph*),leftmargin=1.6em,itemsep=1pt,topsep=2pt]
  \item Có tất cả bao nhiêu cách ngồi thoả mãn các điều kiện trên?
  \item Trong các khẳng định sau, khẳng định nào \emph{chắc chắn đúng}?\\
  (1) Linh ngồi ở một đầu hàng ghế. \quad (2) Hà không ngồi ở đầu hàng ghế. \quad (3) Khoa ngồi ở ghế số $3$.
  \item Nếu biết thêm Linh ngồi ghế số $1$ và Hà ngồi bên trái Khoa, hãy xác định chỗ ngồi của từng bạn.
\end{enumerate}
\end{probbox}
\probsol{
\textbf{Phân tích đề.} Bắt đầu từ điều kiện ``chặt'' nhất (Tùng chỉ có $2$ chỗ), rồi xét $4$ ghế còn lại.
\step{1}{Xét trường hợp Tùng ngồi ghế 5}
Bốn bạn còn lại ngồi ghế $1, 2, 3, 4$. Cặp Hà -- Khoa chiếm hai ghế liền nhau: ghế $(1,2)$, $(2,3)$ hoặc $(3,4)$.
\begin{itemize}[leftmargin=1.4em,itemsep=0pt,topsep=1pt]
  \item Hà -- Khoa ở ghế $(1,2)$: Linh, My phải ngồi ghế $3, 4$ -- cạnh nhau, \textbf{loại}.
  \item Hà -- Khoa ở ghế $(3,4)$: Linh, My ngồi ghế $1, 2$ -- cạnh nhau, \textbf{loại}.
  \item Hà -- Khoa ở ghế $(2,3)$: Linh, My ngồi ghế $1$ và $4$ -- không cạnh nhau, \textbf{được}.
\end{itemize}
Trong trường hợp được: Hà, Khoa đổi chỗ cho nhau được ($2$ cách), Linh, My đổi chỗ cho nhau được ($2$ cách):
$2\times 2 = 4$ cách.
\step{2}{Trường hợp Tùng ngồi ghế 1}
Hoàn toàn tương tự (lật ngược hàng ghế): Hà -- Khoa ở ghế $(3,4)$, Linh và My ở ghế $2$ và $5$: cũng $4$ cách.
\step{3}{Tổng số cách}
$4 + 4 = 8$ cách.
\begin{recallbox}\textbf{Quy tắc nhân:} nếu một việc gồm các bước liên tiếp, bước 1 có $m$ cách, bước 2 có $n$ cách, thì cả việc có
$m\times n$ cách. \textbf{Quy tắc cộng:} các trường hợp tách rời nhau thì cộng số cách.\end{recallbox}
\step{4}{Câu b): kiểm tra từng khẳng định}
\begin{center}\small
\begin{tabular}{|c|c|c|c|c|}
\hline
\rowcolor{tmBlue!10} Ghế 1 & Ghế 2 & Ghế 3 & Ghế 4 & Ghế 5\\ \hline
Linh/My & Hà/Khoa & Hà/Khoa & My/Linh & Tùng\\ \hline
Tùng & Linh/My & Hà/Khoa & Hà/Khoa & My/Linh\\ \hline
\end{tabular}
\end{center}
(1) Sai: khi Tùng ngồi ghế 1, Linh có thể ngồi ghế 2. \ (2) \textbf{Đúng}: trong cả hai dòng, Hà và Khoa luôn ở các ghế giữa
($2, 3$ hoặc $3, 4$). \ (3) Sai: Khoa có thể ngồi ghế $2$ hoặc ghế $4$.
\step{5}{Câu c)}
Linh ngồi ghế $1$ nên Tùng không ở ghế $1$, vậy Tùng ở ghế $5$ (dòng thứ nhất của bảng). Khi đó My ở ghế $4$, Hà -- Khoa ở ghế $2, 3$;
Hà bên trái Khoa nên Hà ghế $2$, Khoa ghế $3$.
\verify{Linh(1), Hà(2), Khoa(3), My(4), Tùng(5): Hà -- Khoa cạnh nhau; Linh ghế 1 và My ghế 4 không cạnh nhau; Tùng ở đầu hàng.
Đúng.}
\ans{a) $8$ cách;\quad b) khẳng định (2);\quad c) Linh -- Hà -- Khoa -- My -- Tùng (ghế 1 đến 5)}
}

% ---------------------------------------------------------------- Bài 35
\begin{probbox}{35}{\relax}
Đánh số các chữ cái $A, B, C, \dots, Z$ (bảng chữ cái tiếng Anh, $26$ chữ) lần lượt là $0, 1, 2, \dots, 25$. \emph{Mật mã Caesar
với khoá $k$} thay mỗi chữ cái mang số $x$ bằng chữ cái mang số là \emph{số dư} của $x + k$ khi chia cho $26$
(dấu cách giữ nguyên).
\begin{figC}
\foreach \L [count=\i from 0] in {A,B,C,D,E,F,G,H,I,J,K,L,M,N,O,P,Q,R,S,T,U,V,W,X,Y,Z}{
  \node[font=\sffamily\small\bfseries,text=Ink] at ({90-\i*360/26}:2.15){\L};
  \node[font=\sffamily\tiny,text=tmBlue] at ({90-\i*360/26}:1.65){\i};}
\draw[Line,line width=0.8pt] (0,0) circle (1.9);
\draw[->,tmOrange,line width=1pt] (110:1.2) arc (110:30:1.2);
\node[font=\scriptsize,text=tmOrange!80!black,align=center] at (0,0){đi theo\\chiều kim\\đồng hồ};
\end{figC}
\begin{enumerate}[label=\alph*),leftmargin=1.6em,itemsep=1pt,topsep=2pt]
  \item Mã hoá dòng chữ \textsf{TOAN HOC} với khoá $k = 5$.
  \item Giải mã dòng chữ \textsf{OVJ CBP}, biết khoá là $k = 7$.
  \item Một tin nhắn có từ \textsf{BAN} được mã hoá thành \textsf{QPC}. Tìm khoá $k$ (với $0 \le k \le 25$), rồi giải mã từ
  \textsf{HPJ} của cùng tin nhắn.
  \item Mã hoá một tin với khoá $10$, rồi mã hoá kết quả thêm một lần nữa với khoá $20$. Hai lần mã hoá đó giống với
  \emph{một} lần mã hoá với khoá nào?
  \item Bạn Tú mã hoá một từ với khoá $4$, rồi lại mã hoá kết quả với khoá $4$, cứ thế nhiều lần. Sau ít nhất bao nhiêu lần
  thì Tú lại thu được đúng từ ban đầu?
\end{enumerate}
\end{probbox}
\probsol{
\textbf{Phân tích đề.} Khi cộng khoá mà vượt quá $25$, ta ``quay vòng'' về đầu bảng chữ cái -- đó chính là lấy \emph{số dư khi
chia cho $26$}. Giải mã là làm ngược lại: \emph{trừ} khoá (nếu ra số âm thì cộng thêm $26$).
\begin{recallbox}Số dư của $a$ khi chia cho $26$ là số $r$ với $0 \le r \le 25$ sao cho $a - r$ chia hết cho $26$.
Ví dụ: $30$ chia $26$ dư $4$; \ $-6$ ``chia $26$ dư'' $20$ vì $-6 + 26 = 20$.\end{recallbox}
\step{1}{Câu a): mã hoá \textsf{TOAN HOC} với $k = 5$}
\begin{center}\small
\begin{tabular}{|c|c|c|c|c|c|c|c|}
\hline
\rowcolor{tmBlue!10} Chữ gốc & T & O & A & N & H & O & C\\ \hline
Số $x$ & $19$ & $14$ & $0$ & $13$ & $7$ & $14$ & $2$\\ \hline
$x + 5$ & $24$ & $19$ & $5$ & $18$ & $12$ & $19$ & $7$\\ \hline
Chữ mã & Y & T & F & S & M & T & H\\ \hline
\end{tabular}
\end{center}
Kết quả: \textsf{YTFS MTH}. (Ở đây không có số nào vượt $25$ nên chưa cần quay vòng.)
\step{2}{Câu b): giải mã \textsf{OVJ CBP} với $k = 7$ (trừ $7$)}
\begin{center}\small
\begin{tabular}{|c|c|c|c|c|c|c|}
\hline
\rowcolor{tmBlue!10} Chữ mã & O & V & J & C & B & P\\ \hline
Số & $14$ & $21$ & $9$ & $2$ & $1$ & $15$\\ \hline
Số $- 7$ & $7$ & $14$ & $2$ & $-5 \to 21$ & $-6 \to 20$ & $8$\\ \hline
Chữ gốc & H & O & C & V & U & I\\ \hline
\end{tabular}
\end{center}
Tin gốc: \textsf{HOC VUI} (Học vui). Chú ý: $2 - 7 = -5$, cộng $26$ được $21$ (chữ V).
\step{3}{Câu c): tìm khoá từ cặp \textsf{BAN} $\to$ \textsf{QPC}}
B ($1$) thành Q ($16$), nên $k = 16 - 1 = 15$. Kiểm tra: A ($0$) $\to 0 + 15 = 15$ (P) đúng; N ($13$) $\to 13 + 15 = 28$,
chia $26$ dư $2$ (C) đúng. Vậy $k = 15$.\\
Giải mã \textsf{HPJ} (trừ $15$): H $= 7$, $7 - 15 = -8 \to 18$ (S); \ P $= 15 \to 0$ (A); \ J $= 9$, $9 - 15 = -6 \to 20$ (U).
Từ gốc là \textsf{SAU}.
\step{4}{Câu d): hai lần mã hoá liên tiếp}
Mỗi chữ được cộng $10$ rồi cộng $20$, tức là cộng $30$. Vì $30$ chia $26$ dư $4$, nên ``cộng $30$ rồi lấy dư'' giống hệt
``cộng $4$ rồi lấy dư''. Hai lần mã hoá giống \textbf{một lần mã hoá với khoá $4$}.
\step{5}{Câu e): mã hoá lặp lại với khoá $4$}
Sau $n$ lần, mỗi chữ đã bị dịch đi $4n$ vị trí. Muốn quay về đúng chữ ban đầu thì phải đi trọn một số vòng, tức $4n$ chia hết
cho $26$. Thử $n = 1, 2, 3, \dots$: \ $4n = 4, 8, 12, \dots$ Số $4n$ nhỏ nhất (khác $0$) chia hết cho $26$ là
$\text{BCNN}(4, 26) = 52$, ứng với $n = 52 : 4 = 13$.
\begin{recallbox}$4n$ chia hết cho $26 = 2\times 13$ $\Leftrightarrow$ $2n$ chia hết cho $13$ $\Leftrightarrow$ $n$ chia hết cho $13$
(vì $13$ là số nguyên tố không chia hết $2$).\end{recallbox}
\begin{tipbox}Với khoá $13$ thì chỉ cần mã hoá $2$ lần là về chỗ cũ ($2\times 13 = 26$). Mã hoá với khoá $13$ hai lần cũng chính là
giải mã!\end{tipbox}
\verify{Câu a): giải mã \textsf{YTFS MTH} bằng cách trừ $5$ được lại \textsf{TOAN HOC}. Câu c): \textsf{SAU} mã hoá với khoá
$15$: S $18 + 15 = 33 \to 7$ (H), A $\to 15$ (P), U $20 + 15 = 35 \to 9$ (J). Đúng \textsf{HPJ}.}
\ans{a) \textsf{YTFS MTH};\ b) \textsf{HOC VUI};\ c) $k = 15$, \textsf{SAU};\ d) khoá $4$;\ e) $13$ lần}
}

% ---------------------------------------------------------------- Bài 36
\begin{probbox}{36}{\stars{1}}
Vẫn đánh số $A = 0, B = 1, \dots, Z = 25$. Một \emph{mật mã nhân} thay chữ cái mang số $x$ bằng chữ cái mang số là
số dư của $3x + 2$ khi chia cho $26$.
\begin{enumerate}[label=\alph*),leftmargin=1.6em,itemsep=1pt,topsep=2pt]
  \item Mã hoá từ \textsf{MEO}.
  \item Giải mã từ \textsf{ECS}.
  \item Bạn Nhật đề nghị dùng $2x + 2$ thay cho $3x + 2$. Hãy chỉ ra hai chữ cái khác nhau có cùng một chữ mã. Vì sao như vậy thì
  không giải mã được?
  \item Giải thích vì sao với $3x + 2$, hai chữ cái khác nhau luôn có hai chữ mã khác nhau.
\end{enumerate}
\end{probbox}
\probsol{
\textbf{Phân tích đề.} Mật mã này không chỉ ``dịch'' mà còn ``nhân'' rồi lấy số dư. Mã hoá thì dễ (tính thẳng); giải mã thì cần
tìm $x$ khi biết số dư của $3x + 2$ -- ta dùng cách ``thêm $26$ cho tới khi chia hết cho $3$''.
\step{1}{Câu a): mã hoá \textsf{MEO}}
M $= 12$: $3\cdot 12 + 2 = 38$, chia $26$ dư $12$ $\to$ M.\\
E $= 4$: $3\cdot 4 + 2 = 14$ $\to$ O.\\
O $= 14$: $3\cdot 14 + 2 = 44$, chia $26$ dư $18$ $\to$ S.\\
Kết quả: \textsf{MOS}. (Thú vị: chữ M được mã hoá thành chính nó.)
\step{2}{Câu b): giải mã \textsf{ECS}}
Với chữ mã mang số $y$, ta cần $x$ ($0 \le x \le 25$) sao cho $3x + 2$ chia $26$ dư $y$, tức là $3x = y - 2$, hoặc
$y - 2 + 26$, hoặc $y - 2 + 52$ (thử lần lượt, lấy số chia hết cho $3$).
\begin{center}\small
\begin{tabular}{|c|c|c|c|c|}
\hline
\rowcolor{tmBlue!10} Chữ mã & $y$ & Thử $y - 2,\ y + 24,\ y + 50$ & $3x$ & Chữ gốc\\ \hline
E & $4$ & $2$ (không), $28$ (không), $54$ (chia hết) & $54 \Rightarrow x = 18$ & S\\ \hline
C & $2$ & $0$ (chia hết) & $0 \Rightarrow x = 0$ & A\\ \hline
S & $18$ & $16$ (không), $42$ (chia hết) & $42 \Rightarrow x = 14$ & O\\ \hline
\end{tabular}
\end{center}
Từ gốc là \textsf{SAO}.
\step{3}{Câu c): vì sao $2x + 2$ không dùng được}
A $= 0$: $2\cdot 0 + 2 = 2$ $\to$ C. \quad N $= 13$: $2\cdot 13 + 2 = 28$, chia $26$ dư $2$ $\to$ C.\\
Hai chữ A và N cùng thành chữ C. Khi nhận được chữ C, người giải mã không biết chữ gốc là A hay N, nên không giải mã chắc chắn
được. (Thực ra mọi số $2x + 2$ đều chẵn, nên chỉ có $13$ chữ mã cho $26$ chữ gốc -- chắc chắn có các cặp trùng nhau.)
\step{4}{Câu d): vì sao $3x + 2$ không bị trùng}
\begin{proofbox}Giả sử hai chữ $x$ và $z$ (với $0 \le z < x \le 25$) có cùng chữ mã. Khi đó $3x + 2$ và $3z + 2$ có cùng số dư khi
chia cho $26$, nên hiệu của chúng $3(x - z)$ chia hết cho $26$. Nhưng $3$ và $26$ không có ước chung nào lớn hơn $1$, nên
$x - z$ phải chia hết cho $26$. Mà $0 < x - z \le 25$ -- không thể chia hết cho $26$. Mâu thuẫn. Vậy hai chữ khác nhau luôn có chữ
mã khác nhau.\end{proofbox}
\begin{warnbox}Với $2x + 2$, lập luận trên hỏng ở chỗ: $2(x - z)$ chia hết cho $26$ chỉ cần $x - z$ chia hết cho $13$, ví dụ
$x - z = 13$ (A và N).\end{warnbox}
\verify{Mã hoá lại \textsf{SAO}: S $3\cdot 18 + 2 = 56 \to 56 - 52 = 4$ (E); A $\to 2$ (C); O $\to 44 - 26 = 18$ (S). Được
\textsf{ECS}, đúng.}
\ans{a) \textsf{MOS};\ b) \textsf{SAO};\ c) A và N đều thành C;\ d) xem chứng minh}
}

% ---------------------------------------------------------------- Bài 37
\begin{probbox}{37}{\relax}
Trong một \emph{hình vuông kì diệu} $3\times 3$, tổng các số trên mỗi hàng, mỗi cột và mỗi đường chéo đều bằng nhau (gọi là
\emph{tổng kì diệu} $M$).
\begin{enumerate}[label=\alph*),leftmargin=1.6em,itemsep=1pt,topsep=2pt]
  \item Gọi $c$ là số ở ô chính giữa. Chứng minh rằng $M = 3c$.
  \item Điền nốt các ô trống của hình vuông kì diệu dưới đây.
\end{enumerate}
\begin{figC}
\lgxcdSud{3}{3}{3}{0/0/9,0/1/19,2/2/13}{}
\end{figC}
\end{probbox}
\probsol{
\textbf{Phân tích đề.} Ở câu b) chỉ biết ba số, không có số ở giữa và cũng không biết $M$. Câu a) chính là ``chìa khoá'': nó liên
hệ số ở giữa với $M$.
\step{1}{Câu a): cộng bốn đường đi qua ô giữa}
Đặt tên các ô:
\begin{figC}
\lgxcdSud{3}{3}{3}{0/0/$a$,0/1/$b$,0/2/$d$,1/0/$e$,1/1/$c$,1/2/$f$,2/0/$g$,2/1/$h$,2/2/$i$}{}
\end{figC}
Có bốn đường đi qua ô giữa: hàng giữa $e + c + f$, cột giữa $b + c + h$, hai đường chéo $a + c + i$ và $d + c + g$. Cộng bốn
tổng này:
\[ (e + c + f) + (b + c + h) + (a + c + i) + (d + c + g) = 4M. \]
Vế trái chứa mỗi ô ngoài đúng \emph{một} lần, còn ô giữa $c$ xuất hiện \emph{bốn} lần. Vậy vế trái bằng
(tổng cả $9$ ô) $+\, 3c$.
\begin{proofbox}Tổng cả $9$ ô bằng tổng của ba hàng, tức $3M$. Do đó $3M + 3c = 4M$, suy ra $3c = M$.\end{proofbox}
\step{2}{Câu b): tìm số ở giữa}
Đường chéo chính đi qua $9$, ô giữa $c$ và $13$: \ $9 + c + 13 = M = 3c$. Suy ra $22 = 2c$, $c = 11$ và $M = 33$.
\step{3}{Điền lần lượt các ô (mỗi lần chọn một đường chỉ còn thiếu một ô)}
\begin{itemize}[leftmargin=1.4em,itemsep=0pt,topsep=1pt]
  \item Hàng trên: $33 - 9 - 19 = 5$ (góc trên phải).
  \item Đường chéo phụ: $33 - 5 - 11 = 17$ (góc dưới trái).
  \item Cột trái: $33 - 9 - 17 = 7$; \quad hàng giữa: $33 - 7 - 11 = 15$.
  \item Hàng dưới: $33 - 17 - 13 = 3$.
\end{itemize}
\begin{figC}
\lgxcdSud{3}{3}{3}{0/0/9,0/1/19,2/2/13}{0/2/5,1/0/7,1/1/11,1/2/15,2/0/17,2/1/3}
\end{figC}
\verify{Cột giữa $19 + 11 + 3 = 33$; cột phải $5 + 15 + 13 = 33$; hàng giữa $7 + 11 + 15 = 33$; hai đường chéo $9 + 11 + 13$ và
$5 + 11 + 17$ đều bằng $33$.}
\begin{tipbox}Kết quả $M = 3c$ đúng với \emph{mọi} hình vuông kì diệu $3\times 3$, kể cả khi có số âm hay phân số.\end{tipbox}
\ans{a) xem chứng minh;\quad b) hàng trên $9, 19, 5$; hàng giữa $7, 11, 15$; hàng dưới $17, 3, 13$}
}

% ---------------------------------------------------------------- Bài 38
\begin{probbox}{38}{\relax}
Ba chiếc cân hai đĩa dưới đây đều thăng bằng. Các vật cùng hình dạng thì nặng bằng nhau.
\begin{figC}
\lgxcdBal{0}{2}{2,3}{5}{1,1,1,1,1}
\lgxcdBal{-1.5}{2}{3,3}{4}{1,1,1,1}
\lgxcdBal{-3}{2}{4,1}{2}{2,2}
\node[font=\small\bfseries,text=Slate] at (-3.3,0){(1)};
\node[font=\small\bfseries,text=Slate] at (-3.3,-1.5){(2)};
\node[font=\small\bfseries,text=Slate] at (-3.3,-3){(3)};
\end{figC}
\begin{enumerate}[label=\alph*),leftmargin=1.6em,itemsep=1pt,topsep=2pt]
  \item Mỗi vật \lgxcdI{3}, \lgxcdI{2}, \lgxcdI{4} nặng bằng bao nhiêu vật \lgxcdI{1}?
  \item Cần bao nhiêu vật \lgxcdI{3} để cân bằng với \lgxcdI{4}\,\lgxcdI{4}\,\lgxcdI{2}\,\lgxcdI{2}?
  \item Biết mỗi vật \lgxcdI{1} nặng $12$ g. Ba vật \lgxcdI{4}\,\lgxcdI{3}\,\lgxcdI{2} nặng tất cả bao nhiêu gam?
\end{enumerate}
\end{probbox}
\probsol{
\textbf{Phân tích đề.} Chọn \lgxcdI{1} làm ``đơn vị đo''. Ta tìm cân nào chỉ có \emph{một loại} vật lạ, giải nó trước, rồi thay
vào các cân còn lại.
\begin{recallbox}Cân thăng bằng giống như một đẳng thức. Ta có thể bỏ (hoặc thêm) những vật giống nhau ở cả hai đĩa, hoặc chia
cả hai đĩa thành các phần bằng nhau, cân vẫn thăng bằng.\end{recallbox}
\step{1}{Cân (2): tìm \lgxcdI{3}}
Hai \lgxcdI{3} cân bằng bốn \lgxcdI{1}. Chia đôi mỗi đĩa: một \lgxcdI{3} cân bằng hai \lgxcdI{1}. \quad
\lgxcdI{3} $= 2$ \lgxcdI{1}.
\step{2}{Cân (1): tìm \lgxcdI{2}}
\lgxcdI{2}\,\lgxcdI{3} cân bằng năm \lgxcdI{1}. Thay \lgxcdI{3} bằng hai \lgxcdI{1}: \lgxcdI{2} cùng hai \lgxcdI{1} cân bằng năm
\lgxcdI{1}. Bỏ hai \lgxcdI{1} ở mỗi đĩa: \lgxcdI{2} $= 3$ \lgxcdI{1}.
\step{3}{Cân (3): tìm \lgxcdI{4}}
\lgxcdI{4}\,\lgxcdI{1} cân bằng hai \lgxcdI{2}, tức là $2\times 3 = 6$ \lgxcdI{1}. Bỏ một \lgxcdI{1} ở mỗi đĩa:
\lgxcdI{4} $= 5$ \lgxcdI{1}.
\step{4}{Câu b)}
\lgxcdI{4}\,\lgxcdI{4}\,\lgxcdI{2}\,\lgxcdI{2} nặng bằng $5 + 5 + 3 + 3 = 16$ \lgxcdI{1}. Mỗi \lgxcdI{3} bằng $2$ \lgxcdI{1},
nên cần $16 : 2 = 8$ vật \lgxcdI{3}.
\step{5}{Câu c)}
\lgxcdI{4}\,\lgxcdI{3}\,\lgxcdI{2} nặng bằng $5 + 2 + 3 = 10$ \lgxcdI{1}, tức là $10\times 12 = 120$ (g).
\begin{tipbox}Nếu gọi khối lượng các vật là $t$ (tam giác), $v$ (vuông), $s$ (sao), $o$ (tròn) thì ba cân chính là ba đẳng thức:
$v + t = 5o$; \ $2t = 4o$; \ $s + o = 2v$. Giải lần lượt được $t = 2o$, $v = 3o$, $s = 5o$.\end{tipbox}
\verify{Cân (1): $3 + 2 = 5$. Cân (2): $2 + 2 = 4$. Cân (3): $5 + 1 = 6 = 3 + 3$. Cả ba cân đều thăng bằng.}
\ans{a) \lgxcdI{3} $= 2$, \lgxcdI{2} $= 3$, \lgxcdI{4} $= 5$ vật \lgxcdI{1};\quad b) $8$ vật \lgxcdI{3};\quad c) $120$ g}
}

% ---------------------------------------------------------------- Bài 39
\begin{probbox}{39}{\relax}
Hình ``bánh xe'' dưới đây có $7$ ô tròn: một ô ở tâm và $6$ ô trên vành, tạo thành $3$ ``nan hoa'' thẳng, mỗi nan gồm $3$ ô
(hai ô đối diện trên vành và ô ở tâm). Cần điền các số $1, 2, 3, 4, 5, 6, 7$, mỗi số vào một ô, sao cho tổng ba số trên mỗi nan
hoa đều bằng cùng một số $S$.
\begin{figC}
\foreach \a in {90,150,210}{\draw[Ink,line width=0.9pt] (\a:1.5)--(\a+180:1.5);}
\draw[Line,line width=0.8pt] (0,0) circle (1.5);
\foreach \a in {30,90,...,330}{\draw[Ink,fill=white,line width=0.9pt] (\a:1.5) circle (0.3);}
\draw[Ink,fill=tmBlue!12,line width=0.9pt] (0,0) circle (0.3);
\end{figC}
\begin{enumerate}[label=\alph*),leftmargin=1.6em,itemsep=1pt,topsep=2pt]
  \item Gọi số ở tâm là $c$. Chứng minh rằng $3S = 28 + 2c$.
  \item Suy ra số ở tâm chỉ có thể là $1$, $4$ hoặc $7$. Với mỗi trường hợp, tìm $S$ và chỉ ra một cách điền.
\end{enumerate}
\end{probbox}
\probsol{
\textbf{Phân tích đề.} Đây là kĩ thuật ``đếm hai lần'': cộng tất cả các nan hoa lại, rồi xem mỗi ô được tính mấy lần.
\step{1}{Câu a): cộng ba nan hoa}
Ba nan hoa có tổng $S + S + S = 3S$. Trong phép cộng này, mỗi ô trên vành được tính \emph{một} lần (mỗi ô vành chỉ nằm trên một
nan), còn ô tâm được tính \emph{ba} lần. Vậy
\[ 3S = (1 + 2 + \dots + 7) + 2c = 28 + 2c. \]
\begin{recallbox}$1 + 2 + 3 + 4 + 5 + 6 + 7 = 28$. Ô tâm đã có mặt một lần trong tổng $28$, nên chỉ cần cộng thêm $2c$ (hai lần
còn lại).\end{recallbox}
\step{2}{Câu b): $c$ nhận giá trị nào?}
$3S$ chia hết cho $3$, nên $28 + 2c$ phải chia hết cho $3$. Thử $c = 1, 2, \dots, 7$:
\begin{center}\small
\begin{tabular}{|c|c|c|c|c|c|c|c|}
\hline
\rowcolor{tmBlue!10} $c$ & $1$ & $2$ & $3$ & $4$ & $5$ & $6$ & $7$\\ \hline
$28 + 2c$ & $\mathbf{30}$ & $32$ & $34$ & $\mathbf{36}$ & $38$ & $40$ & $\mathbf{42}$\\ \hline
Chia hết cho $3$? & có & không & không & có & không & không & có\\ \hline
\end{tabular}
\end{center}
Vậy $c \in \{1; 4; 7\}$, tương ứng $S = 10$; $12$; $14$.
\step{3}{Điền cụ thể}
Mỗi nan hoa: hai ô đối diện phải có tổng $S - c$, nên ta chia $6$ số còn lại thành $3$ cặp có tổng bằng nhau.
\begin{itemize}[leftmargin=1.4em,itemsep=0pt,topsep=1pt]
  \item $c = 1$, $S = 10$: các cặp có tổng $9$: $(2, 7), (3, 6), (4, 5)$.
  \item $c = 4$, $S = 12$: các cặp có tổng $8$: $(1, 7), (2, 6), (3, 5)$.
  \item $c = 7$, $S = 14$: các cặp có tổng $7$: $(1, 6), (2, 5), (3, 4)$.
\end{itemize}
\begin{figC}
\foreach \cx/\cc/\p in {0/1/{2,3,4,7,6,5},4/4/{1,2,3,7,6,5},8/7/{1,2,3,6,5,4}}{
 \begin{scope}[xshift=\cx cm,scale=0.7]
  \foreach \a in {90,150,210}{\draw[Ink,line width=0.8pt] (\a:1.5)--(\a+180:1.5);}
  \foreach \v [count=\k from 0] in \p{\draw[Ink,fill=white,line width=0.8pt] (90+\k*60:1.5) circle (0.33);
     \node[font=\sffamily\small\bfseries,text=tmBlue] at (90+\k*60:1.5){\v};}
  \draw[Ink,fill=tmBlue!12,line width=0.8pt] (0,0) circle (0.33); \node[font=\sffamily\small\bfseries] at (0,0){\cc};
 \end{scope}}
\node[font=\small] at (0,-1.55){$S = 10$}; \node[font=\small] at (4,-1.55){$S = 12$}; \node[font=\small] at (8,-1.55){$S = 14$};
\end{figC}
\begin{warnbox}Đừng nghĩ ``số ở giữa phải là số ở giữa dãy (số $4$)''. Cả $1$ và $7$ cũng đặt ở tâm được!\end{warnbox}
\verify{$S = 10$: $1 + 2 + 7 = 1 + 3 + 6 = 1 + 4 + 5 = 10$. \ $S = 12$: $4 + 1 + 7 = 4 + 2 + 6 = 4 + 3 + 5 = 12$. \
$S = 14$: $7 + 1 + 6 = 7 + 2 + 5 = 7 + 3 + 4 = 14$.}
\ans{a) xem chứng minh;\quad b) $c = 1, S = 10$;\ \ $c = 4, S = 12$;\ \ $c = 7, S = 14$ (cách điền như hình)}
}

% ---------------------------------------------------------------- Bài 40
\begin{probbox}{40}{\stars{1}}
Điền các số $1, 2, 3, 4, 5, 6$, mỗi số vào một ô tròn trên các cạnh của tam giác (ba ô ở ba đỉnh, ba ô ở giữa ba cạnh), sao cho
tổng ba số trên mỗi cạnh đều bằng cùng một số $S$.
\begin{figC}
\draw[Ink,line width=0.9pt] (0,2.2)--(-2.2,-1.3)--(2.2,-1.3)--cycle;
\foreach \p in {(0,2.2),(-2.2,-1.3),(2.2,-1.3),(-1.1,0.45),(1.1,0.45),(0,-1.3)}{\draw[Ink,fill=white,line width=0.9pt] \p circle (0.3);}
\end{figC}
\begin{enumerate}[label=\alph*),leftmargin=1.6em,itemsep=1pt,topsep=2pt]
  \item Gọi $V$ là tổng ba số ở ba đỉnh. Chứng minh rằng $3S = 21 + V$. Từ đó tìm mọi giá trị có thể của $S$.
  \item Với mỗi giá trị của $S$, hãy chỉ ra một cách điền.
  \item Khi $S = 10$, ba số ở ba đỉnh là những số nào? Giải thích vì sao không còn khả năng nào khác.
\end{enumerate}
\end{probbox}
\probsol{
\textbf{Phân tích đề.} Giống Bài 39, ta ``đếm hai lần'': mỗi ô ở đỉnh nằm trên \emph{hai} cạnh, mỗi ô giữa cạnh nằm trên
\emph{một} cạnh.
\step{1}{Câu a): cộng ba cạnh}
Tổng ba cạnh là $3S$. Trong đó mỗi số ở giữa cạnh được tính một lần, mỗi số ở đỉnh được tính hai lần. Vậy
\[ 3S = (1 + 2 + 3 + 4 + 5 + 6) + V = 21 + V. \]
\step{2}{Giới hạn của $V$}
$V$ là tổng của ba số khác nhau trong $1, \dots, 6$: nhỏ nhất $1 + 2 + 3 = 6$, lớn nhất $4 + 5 + 6 = 15$. Ngoài ra
$21 + V$ chia hết cho $3$, mà $21$ chia hết cho $3$, nên $V$ chia hết cho $3$:
\[ V \in \{6;\ 9;\ 12;\ 15\} \quad\Rightarrow\quad S = \frac{21 + V}{3} \in \{9;\ 10;\ 11;\ 12\}. \]
\begin{recallbox}Nếu $a$ chia hết cho $3$ và $a + b$ chia hết cho $3$ thì $b$ cũng chia hết cho $3$.\end{recallbox}
\step{3}{Câu b): một cách điền cho mỗi $S$}
Ta chọn ba số ở đỉnh có tổng $V$, rồi lấy $S$ trừ đi tổng hai số ở hai đầu cạnh để được số ở giữa cạnh đó.
\begin{figC}
\foreach \cx/\S/\t/\l/\r/\ml/\mr/\mb in {0/9/1/2/3/6/5/4, 3.6/10/1/3/5/6/4/2, 7.2/11/2/4/6/5/3/1, 10.8/12/4/5/6/3/2/1}{
 \begin{scope}[xshift=\cx cm,scale=0.62]
  \draw[Ink,line width=0.8pt] (0,2.2)--(-2.2,-1.3)--(2.2,-1.3)--cycle;
  \foreach \p/\v/\col in {(0,2.2)/\t/Ink,(-2.2,-1.3)/\l/Ink,(2.2,-1.3)/\r/Ink,(-1.1,0.45)/\ml/tmBlue,(1.1,0.45)/\mr/tmBlue,(0,-1.3)/\mb/tmBlue}{
    \draw[Ink,fill=white,line width=0.8pt] \p circle (0.36); \node[font=\sffamily\small\bfseries,text=\col] at \p {\v};}
  \node[font=\small] at (0,-2.2){$S = \S$};
 \end{scope}}
\end{figC}
Ví dụ với $S = 9$: đỉnh $1, 2, 3$; cạnh $1$--$2$ cần $9 - 3 = 6$; cạnh $1$--$3$ cần $9 - 4 = 5$; cạnh $2$--$3$ cần $9 - 5 = 4$.
\step{4}{Câu c): $S = 10$ thì $V = 3\cdot 10 - 21 = 9$}
Các bộ ba số khác nhau có tổng $9$: $\{1, 2, 6\}$, $\{1, 3, 5\}$, $\{2, 3, 4\}$. Thử từng bộ:
\begin{itemize}[leftmargin=1.4em,itemsep=0pt,topsep=1pt]
  \item Đỉnh $\{1, 2, 6\}$: cạnh nối $1$ và $2$ cần ô giữa là $10 - 3 = 7$ -- không có số $7$. \textbf{Loại.}
  \item Đỉnh $\{2, 3, 4\}$: cạnh nối $3$ và $4$ cần $10 - 7 = 3$ -- nhưng $3$ đã ở đỉnh. \textbf{Loại.}
  \item Đỉnh $\{1, 3, 5\}$: các ô giữa là $10 - 4 = 6$, $10 - 6 = 4$, $10 - 8 = 2$ -- đúng là ba số còn lại $2, 4, 6$.
  \textbf{Được.}
\end{itemize}
Vậy khi $S = 10$, ba số ở đỉnh bắt buộc là $1, 3, 5$.
\begin{warnbox}Điều kiện ``$V$ chia hết cho $3$'' chỉ là điều kiện \emph{cần}: nó cho biết $S$ \emph{có thể} bằng bao nhiêu, nhưng
vẫn phải chỉ ra cách điền cụ thể (câu b) và kiểm tra từng bộ (câu c).\end{warnbox}
\begin{parentbox}Cho con cắt $6$ mẩu giấy ghi số $1$ đến $6$ và xếp thử trên hình vẽ. Sau vài lần thử, hỏi con: ``Số ở đỉnh được
cộng mấy lần?'' -- đó là ý chính của lời giải.\end{parentbox}
\verify{$S = 10$: $1 + 6 + 3 = 3 + 4 + 5 = 5 + 2 + 1 = 10$. $S = 12$: $4 + 3 + 5 = 5 + 1 + 6 = 6 + 2 + 4 = 12$.}
\ans{a) $S \in \{9; 10; 11; 12\}$;\quad b) như hình;\quad c) đỉnh là $1, 3, 5$}
}
